\section{Split Protocols and Identifier Probes}
\label{sec:protocols}

\subsection{Protocols}

We implement six partitions and report four. \textbf{Random flow} is stratified
random assignment of individual flows, the protocol we take to be standard in
this literature; we report it only as a contrast. \textbf{Session-disjoint}
assigns whole capture sessions --- a capture file for 5G Traffic, a capture day
for CESNET --- to exactly one side. \textbf{Temporal} sorts by first-packet
timestamp and takes the earliest flows for training. \textbf{Server-disjoint}
assigns whole server addresses to one side. Every protocol is materialised as a
committed manifest carrying per-side row indices and a SHA256 over the flow
identifiers, so a reader can verify that a reported number was produced on the
released partition.

\subsection{Identifier probes}

Before training anything, we fit a depth-limited decision tree on a
\emph{single} identifier column and score it under each protocol. The
identifiers are destination port, SNI string, server address and \texttt{/24}
prefix, destination AS, capture-session identifier, and transport protocol. A
probe whose column is absent or constant in a corpus reports
\emph{unavailable}, never a degenerate number --- a dataset with no SNI must
not be reported as one where SNI does not help.

% Generated by scripts/make_paper_assets.py. Do not edit by hand.
\begin{table}[t]
\centering
\small
\begin{tabular}{lcc}
\toprule
Split protocol & cesnet\_quic22 & fiveg\_traffic \\
\midrule
\multicolumn{3}{l}{\emph{FlowCon-X (prototype head)}} \\
Random flow & 0.781 $\pm$ 0.006 & 0.991 $\pm$ 0.000 \\
Session-disjoint & \textsc{todo} & 0.522 $\pm$ 0.139 \\
Temporal & 0.781 $\pm$ 0.005 & \textsc{todo} \\
Server-disjoint & 0.588 $\pm$ 0.007 & \textsc{todo} \\
\midrule
\multicolumn{3}{l}{\emph{AppScanner (2016), identical splits}} \\
Random flow & 0.768 & 0.996 \\
Session-disjoint & 0.772 & 0.640 \\
Temporal & 0.761 & 0.255 \\
Server-disjoint & 0.566 & 0.711 \\
\bottomrule
\end{tabular}
\caption{Macro-F1 under four split protocols, every method on the identical committed splits. The random-flow row is the protocol standard in the literature and is reported only as a contrast: correlated flows from one capture session appear on both sides of it. \textbf{The size of that contrast is dataset-dependent}, which is itself the result: on a corpus of single-session captures a flow's capture is nearly its label, while on a backbone corpus whose flows are already independent the protocol makes little difference. The caption deliberately does not assert a direction; read the columns.}
\label{tab:split-contrast}
\end{table}


\subsection{What the probes find}

\paragraph{5G Traffic: the capture is the label.} Under a random split the
server address alone reaches macro-F1 0.999 and the capture-session identifier
alone 0.983. A classifier that learned nothing except which recording a flow
came from would score in the range published work reports for this corpus.
Under session-disjoint splitting the capture probe falls to \textbf{exactly
0.000} --- the grouping doing precisely what it claims --- and AppScanner falls
from 0.996 to 0.640, our model from 0.992 to 0.567. Temporal splitting is
harder still: 0.255 and 0.391 respectively. These captures span five months, so
a temporal split separates application versions and network conditions as well
as time, and it is the protocol under which this corpus behaves least like the
one the literature reports.

\paragraph{CESNET-QUIC22: the server is the label, and only the server.} We
evaluated five protocols at eight seeds each and tested each against
session-disjoint with a Wilcoxon signed-rank test, Holm-Bonferroni corrected
across the family:

\begin{center}
\small
\begin{tabular}{lrrr}
\toprule
Split protocol & Macro-F1 & Cohen's $d$ & $p$ \\
\midrule
Session-disjoint (ref.) & $0.7798 \pm 0.0061$ & --- & --- \\
Random flow & $0.7805 \pm 0.0057$ & $-0.12$ & 0.742 \\
Temporal & $0.7811 \pm 0.0051$ & $-0.31$ & 0.547 \\
Client-disjoint & $0.7791 \pm 0.0055$ & $0.13$ & 0.945 \\
\textbf{Server-disjoint} & $\mathbf{0.5881 \pm 0.0068}$ & $\mathbf{17.03}$ & $\mathbf{0.0078}^{*}$ \\
\bottomrule
\end{tabular}
\end{center}

\textbf{Four of five axes make no difference; the fifth costs 0.192 macro-F1.}
The three nulls are properly powered rather than absences of evidence: eight
seeds, effect sizes below 0.35, $p > 0.5$. Only server-disjointness survives
the correction ($^{*}$).

The client-disjoint row is worth its own sentence, because we added it
expecting the opposite. A third of the 48{,}072 distinct client addresses in
this corpus appear on more than one capture day, so grouping by day
demonstrably does not separate clients --- and separating them changes nothing
($p = 0.945$). Backbone flows from one client on different days are no more
alike, for this task, than flows from different clients.

AppScanner behaves the same way: 0.768, 0.772 and 0.762 under random,
session-disjoint and temporal splitting. Server-disjoint splitting takes it to 0.566 and our model from 0.781 to 0.588,
and takes the server-address probe from 0.772 to \textbf{0.026} --- the floor,
confirming the partition does what it says. Table~\ref{tab:split-contrast}
shows the two corpora side by side: the CESNET column is flat across three
protocols and drops only on the fourth, while the 5G column drops on all
three.

\paragraph{The labelling ceiling.} The SNI probe reaches 0.967 on CESNET under
every protocol except server-disjoint, where it still reaches 0.791. This is
expected, since that corpus's labels are derived from SNI, and it is the most
important number in this section: no result on CESNET-QUIC22 can be read as
matching ground truth, because the ground truth \emph{is} a name-based
heuristic and a decision tree recovers it at 0.967.

\subsection{The rule we propose}

The two corpora fail on different axes, and neither protocol protects against
the other's failure mode. We propose that the axis be chosen by measurement
rather than convention:

\begin{quote}
Before reporting any classifier result, fit a single-identifier probe for each
identifier the corpus retains, under each candidate protocol. Adopt the
protocol under which every identifier probe collapses toward the majority-class
floor. Report the probe table alongside the results.
\end{quote}

This costs seconds and it is falsifiable: if no protocol collapses the probes,
the corpus cannot support the claim being made of it. On 5G Traffic the rule
selects session-disjoint or temporal; on CESNET-QUIC22 it selects
server-disjoint. Adopting session-disjoint splitting on CESNET --- the choice a
practitioner following current best practice would make --- controls nothing
there, and our measurements say so.
